\begin{abstract}
Constructive and destructive gradient regimes coexist with learning progress: training loss can decrease even when consecutive updates repeatedly undo one another. We characterize these regimes through the alignment of consecutive gradients and trace their effects from parameter updates to changes in the learned function. For smooth full-batch gradient descent, we show that consecutive gradients recover the directional curvature averaged over the completed step. This diagnostic tracks the transition to oscillatory training without additional derivative evaluations. We then trace hidden representations and predictions on fixed inputs around Adam loss spikes in an MLP, a ResNet, and a Transformer. In all three spike windows, predictions oscillate predominantly within a low-dimensional functional subspace, leaving a net output displacement of only 1.7--3.5\% of the distance traveled. In ResNet and the Transformer, predictions can reverse while earlier representations continue moving constructively. Motivated by these observations, we cap the direct change that Adam's proposed update would produce in individual hidden units. In matched continuations across all three models, the cap dampens loss spikes and lowers mean training loss while functional reversals persist. These findings point toward interventions that go beyond loss-landscape geometry by using the network's internal responses to guide parameter updates.
\end{abstract}
